\chapter{Syntax of {\meq} Configuration Files}
\label{ap:config}

A run of the {\meq} program is described entirely by one TOML file, and this
appendix is the complete reference for it. \Chref{ch:using} is the narrative
version.

\section{Shape of the file}

The file is a table of tables. Exactly eight top-level tables are part of the
schema:

\begin{center}
\begin{tabular}{l l}
	\hline
	Table & \\
	\hline
	\texttt{[mesh]} & \textbf{required} \\
	\texttt{[discretisation]} & \textbf{required} \\
	\texttt{[source]} & \textbf{required} \\
	\texttt{[boundary]} & optional; contains \texttt{[boundary.shape]} \\
	\texttt{[solver]} & optional \\
	\texttt{[output]} & optional \\
	\texttt{[initialguess]} & optional \\
	\texttt{[adaptivity]} & optional \\
	\hline
\end{tabular}
\end{center}

An absent optional table behaves exactly like a present but empty one: every key
falls back to its default. Note the spellings \texttt{[initialguess]} and
\texttt{[adaptivity]} --- neither carries a separator, and the second is not
\texttt{[adaptive]}.

\begin{quote}
	\textbf{Table names are lower case and keys are \texttt{UpperCamelCase}.}
	This is a deliberate mismatch with the C++ naming convention used
	everywhere else in the tree, inherited from a sibling project, and it is
	not an oversight to be reconciled. The two nested table \emph{keys},
	\texttt{shape} and \texttt{species}, follow the table rule rather than the
	key rule: they are written \texttt{[boundary.shape]} and
	\texttt{[[source.species]]}.
\end{quote}

\section{How the parser behaves}

\subsection{Everything is validated before anything is done}

Construction of the configuration either succeeds and leaves every accessor
meaningful, or throws. There is no partially valid configuration, and no work of
any kind --- no mesh, no source, no solve --- begins until the whole file has
been accepted.

\subsection{Unknown keys and tables are errors}

Not warnings, and not ignored. The message names the nearest accepted spelling
when there is one within a small edit distance, and lists the keys the table
does accept.

This is the entire point of having a schema. A key that is silently ignored is a
run that quietly did something other than what the file says, and no amount of
care in reading the output recovers from that.

\subsection{Types}

\begin{description}

\item[Numbers] accept either TOML spelling: \texttt{RMin = 0} and
	\texttt{RMin = 0.0} are both read as the float zero.

\item[Counts] do not. \texttt{NR = 4.0} is refused rather than truncated.

\item[Booleans] must be real booleans. \texttt{Enabled = "true"} is an error,
	not a silent false.

\item[Arrays] are validated element by element, and the diagnostic names the
	element.

\end{description}

That first distinction is worth more than it appears, and it cost a real bug
once. Reading every number as a float rejects \texttt{RMin = 0}; but the natural
remedy in the TOML library used here returns the \emph{default} when a
conversion fails, because a failed conversion is indistinguishable from a missing
key. That turns a loud failure into a silent wrong answer in the configuration.
{\meq} therefore checks each node's type explicitly.

\subsection{Diagnostics}

Every fault is a \texttt{meq::ConfigError} carrying the file name and the fully
qualified key, including array elements --- so a fault in the third species
block is reported against \texttt{source.species[2].Mass} and not merely against
\texttt{source}.

\subsection{Paths}

Profile and mesh paths are resolved against the \emph{working directory of the
run}, not against the directory the configuration file lives in.

\section{\texttt{[mesh]} --- required}

The background mesh: either a box that {\meq} triangulates, or a file.

\begin{description}

\item[\texttt{File}] \emph{string, default \texttt{""}.} A mesh in any format
	MFEM reads. A non-empty value selects file mode, in which the box keys are
	accepted but not used.

\item[\texttt{RMin}] \emph{float, required unless \texttt{File} is set.} Lower
	radial bound, in metres. Must not be negative: $r$ is a cylindrical radius,
	and a box reaching $r = 0$ contains the operator's singularity.

\item[\texttt{RMax}] \emph{float, required unless \texttt{File} is set.} Must
	exceed \texttt{RMin}.

\item[\texttt{ZMin}, \texttt{ZMax}] \emph{float, required unless \texttt{File}
	is set.} Vertical bounds; \texttt{ZMax} must exceed \texttt{ZMin}.

\item[\texttt{NR}, \texttt{NZ}] \emph{integer, default 1.} Cells across each
	direction \emph{before} refinement; each must be at least~1. Cells are split
	into triangles, so the initial element diameter is the cell diagonal.

\item[\texttt{RefinementLevels}] \emph{integer, default 0.} Levels of uniform
	refinement, each halving $h$. Applies whether the mesh came from a box or a
	file. Must not be negative.

\end{description}

\begin{quote}
	\textbf{\texttt{File} is effectively supported on the fitted path only.} The
	box bounds are what the gridded output samples over and what sets the search
	length for the curved boundary's transfer paths. With a mesh file and no
	\texttt{[boundary.shape]} to supply a bounding box, both are degenerate.
\end{quote}

\section{\texttt{[discretisation]} --- required}

\begin{description}

\item[\texttt{PolynomialDegree}] \emph{integer, required.} The degree $k$ of all
	three HDG spaces. Must not be negative. The source papers report $k = 1$
	to~5; beyond about~5, or eight refinements, round-off dominates.

\item[\texttt{Tau}] \emph{float, default 1.0.} The HDG stabilisation,
	dimensionless and positive. Optimal order needs only $\tau = O(1)$. Read
	\chref{ch:using} before changing it.

\end{description}

\section{\texttt{[source]} --- required}

\texttt{Type} is required and discriminates. The accepted values are exactly
\texttt{"soloviev"}, \texttt{"mhd"}, \texttt{"manufactured"} and
\texttt{"rotating"}.

\begin{quote}
	\textbf{The accepted key set depends on \texttt{Type}.} A key belonging to
	another source is rejected as unknown, not ignored --- so
	\texttt{Normalised} under a Solov'ev source is an error.
\end{quote}

\subsection{\texttt{Type = "soloviev"}}

$F = -\left((1-A)r^{2} + A\right)$, independent of $\psi$. The problem is linear
and Newton converges in one step.

\begin{description}
\item[\texttt{A}] \emph{float, required.} Dimensionless, with the flux
	normalised so that $A + C = 1$. The NSTX benchmark uses \texttt{-0.52}.
\end{description}

\subsection{\texttt{Type = "manufactured"}}

The nonlinear manufactured solution $\psi = \sin(K_r(r + R_0))\cos(K_z z)$ of
S\'anchez-Vizuet and Solano \cite{SanchezVizuetSolano2019}, whose
$\psi$-dependence is linear, quadratic and exponential at once.

\begin{description}
\item[\texttt{R0}] \emph{float, required.} The radial \emph{offset} inside the
	sine. \textbf{Not a major radius.}
\item[\texttt{Kr}, \texttt{Kz}] \emph{float, required.} Radial and vertical
	wavenumbers.
\end{description}

\subsection{\texttt{Type = "mhd"}}

$F = \mu_{0}\,r^{2}p'(\psi) + (gg')(\psi)$, from two tabulated profiles. Both
are files: there is no scalar spelling on this path.

\begin{description}

\item[\texttt{PPrimeFile}] \emph{string, required, non-empty.} Path to a
	tabulated $\mathrm{d}p/\mathrm{d}\psi$. The file format is
	\S\ref{sec:profileformat}.

\item[\texttt{GGPrimeFile}] \emph{string, required, non-empty.} Path to a
	tabulated $g\,\mathrm{d}g/\mathrm{d}\psi$ --- what EQDSK calls
	\texttt{FF'}.

\item[\texttt{PPrimeScale}, \texttt{GGPrimeScale}] \emph{float, default 1.0,
	finite.} A constant multiplying the table as read --- value, slope and
	second derivative alike, so a scaled table stays self-consistent. This is
	the units mechanism: a file stays readable in its own units and the
	conversion goes in the scale. The abscissa is never scaled.

\item[\texttt{Mu0}] \emph{float, default the SI $\mu_0$, positive.} The vacuum
	permeability multiplying \emph{only} the $r^{2}p'$ term. Set it to 1.0 for a
	run in normalised units.

\item[\texttt{Normalised}] \emph{boolean, default false.} The tables are
	functions of $\Psi = \psi/\psiax$ rather than of $\psi$. This makes $\psiax$
	an unknown of the nonlinear system, closed by a bordered Newton.

\item[\texttt{PsiAxis}] \emph{float, required when \texttt{Normalised} is true
	and refused otherwise.} A starting \emph{guess} for the Newton iteration,
	not a scale factor. Must be finite and non-zero, since $\Psi = \psi/\psiax$
	is undefined at zero.

\item[\texttt{ProfileFile}] is reserved for reading several profiles out of one
	NetCDF file and is always refused, with a message saying so.

\item[\texttt{PPrimeVariable}, \texttt{PPrimeFit}, \texttt{GGPrimeVariable},
	\texttt{GGPrimeFit}] are reserved for the same unwritten facility. They are
	in this path's accepted-key list and are inert; on the rotating path they
	are refused outright.

\end{description}

\subsection{\texttt{Type = "rotating"}}

A rotating multi-species plasma: the generalised Grad--Shafranov equation of
Abel and co-workers \cite{Abel2013}, closed by quasineutrality, with
$p = \sum_s n_s T_s$.

\begin{description}

\item[\texttt{species}] \emph{array of tables, required.} Written
	\texttt{[[source.species]]}; see \S\ref{sec:species}. Between two and eight
	of them --- fewer than two and quasineutrality has nothing to solve; the
	upper bound exists so that the per-quadrature-point work allocates nothing.

\item[\texttt{Omega} \emph{or} \texttt{OmegaFile}] \emph{float or string,
	default: neither.} The rigid rotation frequency $\omega(\psi)$, in radians
	per second, as a constant or as a table. At most one of the two.
	\textbf{Giving neither means no rotation}, and the source reduces to the
	static equation.

\item[\texttt{GGPrime} \emph{or} \texttt{GGPrimeFile}] \emph{float or string;
	exactly one is required.} $g\,\mathrm{d}g/\mathrm{d}\psi$, as a constant or
	as a table.

\item[\texttt{OmegaScale}, \texttt{GGPrimeScale}] \emph{float, default 1.0,
	finite.} A scale on whichever form was given. A scale with nothing to scale
	is an error.

\item[\texttt{ReferenceRadius}] \emph{float, default 1.0, positive.}
	\textbf{The gauge.} The radius at which the electrostatic potential
	vanishes, and therefore the radius at which each species' \texttt{Density}
	is the physical density. It is a constant radius --- not the magnetic axis,
	and not a flux-surface average.

\item[\texttt{Mu0}, \texttt{Normalised}, \texttt{PsiAxis}] as for
	\texttt{"mhd"}, with the same conditions and the same messages.

\item[Reserved keys] Five keys are reserved for reading several profiles out
	of one NetCDF file, which is not implemented, and all five are
	\emph{refused} rather than ignored on this path:
	\texttt{ProfileFile}, \texttt{OmegaVariable}, \texttt{OmegaFit},
	\texttt{GGPrimeVariable} and \texttt{GGPrimeFit}.

\end{description}

\subsection{\texttt{[[source.species]]}}
\label{sec:species}

One block per species, under \texttt{[source]}, for the rotating source only.

\begin{description}

\item[\texttt{Name}] \emph{string, default \texttt{"species}$i$\texttt{"} with
	$i$ counted from zero.} Used in output variable names, sanitised to
	something a NetCDF name may be.

\item[\texttt{Mass}] \emph{float, required, positive.} Particle mass in
	kilogrammes.

\item[\texttt{Charge}] \emph{float, required, non-zero.} $Z_s$: \textbf{signed
	and dimensionless}, so $+1$ for a proton and $-1$ for an electron. Not
	coulombs.

\item[\texttt{Temperature} \emph{or} \texttt{TemperatureFile}] \emph{float or
	string; exactly one is required.} $T_s$ \textbf{in joules}, as a constant or
	as a table.

\item[\texttt{TemperatureScale}] \emph{float, default 1.0, finite.} How a file
	written in electronvolts or kiloelectronvolts stays readable.

\item[\texttt{Density} \emph{or} \texttt{DensityFile}] \emph{float or string;
	exactly one is required unless this species is \texttt{Neutralising}, in
	which case both are refused.} $n_{s0}$ in inverse cubic metres, \emph{on the
	curve} $r = \texttt{ReferenceRadius}$.

\item[\texttt{DensityScale}] \emph{float, default 1.0, finite.} Refused when
	there is no density to scale --- including on the neutralising species.

\item[\texttt{Neutralising}] \emph{boolean, default false.} This species'
	density is \emph{derived} from the others by charge neutrality.
	\textbf{Exactly one species must set it}, because fixing the gauge removes
	one function's worth of freedom from the densities.

\item[Reserved keys] Four more are reserved for the same unwritten facility
	and are likewise refused: \texttt{TemperatureVariable},
	\texttt{TemperatureFit}, \texttt{DensityVariable} and
	\texttt{DensityFit}.

\end{description}

Two further cross-species rules, each with its own message: exactly one species
must be neutralising, and the species must between them carry charges of both
signs or quasineutrality has no solution.

\section{\texttt{[boundary]}}

\begin{description}

\item[\texttt{Type}] \emph{string, default \texttt{"zero"}.} Either
	\texttt{"zero"}, which imposes $\psi = 0$ --- the fixed-boundary problem
	proper --- or \texttt{"exact"}, which imposes the source's known closed-form
	solution and is a convergence-study device.

\end{description}

\begin{quote}
	\textbf{\texttt{Type = "exact"} is refused by the program}, with an
	explanation and exit code~1. A \texttt{meq::Source} carries $F$ and
	$\partial F/\partial\psi$ and no closed form; the closed forms live in the
	test fixtures, where they are the thing being converged against. It remains
	usable from the library, which is why the key exists.
\end{quote}

\subsection{\texttt{[boundary.shape]}}

Present with a real shape, this selects the \emph{curved} path of
\chref{ch:using}.

\begin{description}

\item[\texttt{Type}] \emph{string, default \texttt{"none"}.} One of
	\texttt{"none"}, \texttt{"miller"} or \texttt{"mxh"}. With \texttt{"none"}
	every other key here is accepted and ignored.

\item[\texttt{R0}] \emph{float, required unless \texttt{"none"}.} Major radius
	of the shape's centre, in metres.

\item[\texttt{Z0}] \emph{float, default 0.0.} Centre height.

\item[\texttt{MinorRadius}] \emph{float, required unless \texttt{"none"}.}

\item[\texttt{Elongation}] \emph{float, default 1.0.} $\kappa$.

\item[\texttt{Triangularity}] \emph{float, default 0.0; \texttt{"miller"} only.}
	$\delta$, which \textbf{enters as $\arcsin\delta$} and must therefore
	satisfy $|\delta| < 1$.

\item[\texttt{Squareness}] \emph{float, default 0.0; \texttt{"miller"} only.}
	Zero gives the original three-parameter Miller shape \cite{Miller1998}.

\item[\texttt{CosCoefficients}] \emph{array of floats, default empty;
	\texttt{"mxh"} only.} $c_0, c_1, \ldots$ --- \textbf{starting at $c_0$},
	which is the tilt.

\item[\texttt{SinCoefficients}] \emph{array of floats, default empty;
	\texttt{"mxh"} only.} $s_1, s_2, \ldots$ --- \textbf{starting at $s_1$};
	there is no $s_0$.

\end{description}

Each shape refuses the other's keys rather than ignoring them, and
\texttt{"mxh"} with no harmonics at all is refused with a message suggesting
\texttt{"miller"}, since that is an ellipse.

Four further conditions are checked when the shape is built rather than when the
file is parsed, and each exits~1 with a message naming it: the minor radius, the
major radius and the elongation must be positive; the surface must not reach or
cross the axis, where the operator's $1/r$ is not integrable; and the surface
must be star-shaped about its own centre, since the level set is built by radial
bisection and cannot be built otherwise.

Two more come from the interaction with the mesh. The shape must enclose at
least one background element --- otherwise the mesh is too coarse for it, or it
lies outside the box --- and it must not touch the edge of the box, since then
part of $\Gamma_h$ is inherited from the box, is fitted rather than transferred,
and carries no datum.

\section{\texttt{[solver]}}

\begin{description}
\item[\texttt{NewtonMaxIterations}] \emph{integer, default 20, at least 1.}
\item[\texttt{NewtonRelativeTolerance}] \emph{float, default $10^{-8}$,
	positive.} Relative to the residual at the \emph{cold} iterate rather than
	at the iterate handed in --- see \chref{ch:embedding} for why that
	distinction exists and what it prevents.
\item[\texttt{NewtonAbsoluteTolerance}] \emph{float, default $10^{-12}$,
	positive.}
\end{description}

\begin{quote}
	\textbf{\texttt{LinearMaxIterations} and \texttt{LinearTolerance} are
	refused,} with a message saying why. They configure an \emph{iterative}
	inner solve, and {\meq} solves the hybridized trace system with a
	\emph{direct} one --- UMFPACK, PARDISO or cuDSS --- which has neither an
	iteration count nor a tolerance to set. They were once parsed and validated
	and read by nothing, which is worse than either accepting or rejecting them:
	a key that validates is a key its author believes is doing something. Delete
	the line from any file that carries it.
\end{quote}

\section{\texttt{[output]}}

\begin{description}

\item[\texttt{Directory}] \emph{string, default \texttt{"."}.} \textbf{Not
	created by {\meq}.} A missing directory is exit code~3, after the
	equilibrium has been computed.

\item[\texttt{Prefix}] \emph{string, default \texttt{"meq"}, non-empty.} The
	stem of every output file name.

\item[\texttt{GridNR}, \texttt{GridNZ}] \emph{integer, default 129, at least 2.}
	Sampling \emph{nodes} for the gridded output, so the spacing is
	$(R_{\max} - R_{\min})/(\texttt{GridNR} - 1)$. Nothing whatever to do with
	\texttt{[mesh] NR}.

\end{description}

\section{\texttt{[initialguess]}}

\begin{description}

\item[\texttt{Type}] \emph{string, default \texttt{"none"}.} One of
	\texttt{"none"}, which starts from the Dirichlet datum; \texttt{"ramp"},
	which makes $\psi$ run from $-\texttt{Amplitude}$ to $+\texttt{Amplitude}$
	across $z$; or \texttt{"gridfunction"}, which reads a stored answer.

\item[\texttt{Amplitude}] \emph{float, default 0.3, positive; \texttt{"ramp"}
	only.} The point of the ramp is that $\psi$ crosses zero in the
	\emph{interior}, so an amplitude of zero puts the iteration straight onto
	the trivial branch the ramp exists to avoid.

\item[\texttt{File}, \texttt{MeshFile}] \emph{string, both required for
	\texttt{"gridfunction"}.} The stored grid function and the mesh it lives on
	--- a grid function cannot be read without its mesh.

\end{description}

The ramp is not a nicety: several standard sources vanish at $\psi = 0$, so with
homogeneous data the trivial solution is a genuine solution and Newton lands on
it. A stored guess on the \emph{same} mesh is an exact restart, every
coefficient; on a different one it is interpolated. It is applied on the first
cycle only --- later adaptive cycles warm-start from the previous cycle.

\section{\texttt{[adaptivity]}}

\begin{description}

\item[\texttt{Enabled}] \emph{boolean, default false.}

\item[\texttt{MaxIterations}] \emph{integer, default 10, at least 1.} Counts
	\textbf{solves}, not refinements: 4 means at most three refinements, and 1
	is a single solve with an estimate printed.

\item[\texttt{Strategy}] \emph{string, default \texttt{"doerfler"}.} Either
	\texttt{"doerfler"} or \texttt{"maximum"}, both lower case.

\item[\texttt{Theta}] \emph{float, default 0.6, in $(0, 1]$.} For D\"orfler
	marking, the fraction of the total estimated error the marked elements must
	carry. \textbf{The two strategies read it oppositely}, which is why the
	strategy is named rather than inferred.

\item[\texttt{TargetError}] \emph{float, default $10^{-6}$, positive.} Absolute,
	in the estimator's own norm.

\end{description}

\section{Profile files}
\label{sec:profileformat}

Every \texttt{*File} key names a text table read into a piecewise Hermite cubic.
One knot per line, three whitespace-separated numbers:

\begin{verbatim}
    # a comment
    psi     f(psi)     f'(psi)
\end{verbatim}

\textbf{All three columns are required.} The interpolant passes through the
value \emph{and} the slope, which is what lets the profile supply its own
derivative --- and \chref{ch:intro} explains why a profile that cannot do that
cannot be used at all.

\begin{description}

\item[Comments] A line whose first non-blank character is \texttt{\#}.

\item[Termination] A \emph{blank line ends the table}, so several profiles may
	be concatenated in one stream. The corollary is that there must be no blank
	line between the header comments and the numbers.

\item[Constraints] At least two knots; $\psi$ strictly increasing; every value
	finite.

\item[Outside the table] The profile is extended by a \emph{constant}: the value
	is clamped to the nearest end knot and the first and second derivatives are
	exactly zero. Nothing throws, because a Newton iterate routinely overshoots
	the physical range. The extension is continuous but not smooth.

\item[Second derivatives] A Hermite cubic is only continuously differentiable
	once, so the second derivative is piecewise linear and \emph{jumps at every
	interior knot}. It is used --- a rotating source's Jacobian needs it --- and
	nothing has measured what the jumps cost the Newton iteration.

\end{description}

\begin{quote}
	\textbf{A table written against $\psi$ and a table written against
	$\Psi = \psi/\psiax$ are different files, and neither will complain.}
	Re-expressing $f$ against the normalised flux divides the first column by
	the axis flux \emph{and multiplies the derivative column by it}. Hand one to
	the wrong configuration and it parses, solves and converges --- to a plasma
	whose gradient is out by that factor. Say in the file's header which one it
	is; the shipped examples do.
\end{quote}

\section{The shipped examples}

\texttt{examples/} in the source tree holds a worked instance of each path.

\begin{description}

\item[\texttt{soloviev-nstx.toml}] The headline analytic benchmark: a Solov'ev
	source on a fitted rectangle at high degree, linear in $\psi$ so that Newton
	is one step. It is deliberately \emph{not} the published NSTX equilibrium,
	since $\Gamma$ here is the box rather than the separatrix.

\item[\texttt{manufactured.toml}] The manufactured nonlinear solution with an
	exact boundary condition --- a convergence-study file the program refuses,
	kept because the library path uses it.

\item[\texttt{manufactured-driver.toml}] The same source with a zero boundary
	condition: the Newton path end to end, and runnable.

\item[\texttt{mhd-rectangle.toml}] The tabulated-profile source, genuinely
	semi-linear. The one example whose input a user would actually write.

\item[\texttt{miller-curved.toml}] The curved boundary by transfer, on a
	background rectangle. The configuration {\meq} exists to run.

\item[\texttt{miller-adaptive.toml}] The same shape solved adaptively, with the
	companion-mesh domain update.

\item[\texttt{rotating-rectangle.toml}] Sonic toroidal rotation: two species,
	one of them neutralising, with the reference radius as the gauge.

\item[\texttt{rotating-normalised.toml}] The same plasma with its profiles in
	normalised flux, so that the axis flux is an unknown closed by the bordered
	Newton.

\item[\texttt{*.dat}] The profile tables the two source types above read.
	Between them they carry the fullest statement of the file format anywhere in
	the tree, and the two density files are the same profile written against
	$\psi$ and against $\Psi$ --- the trap of \S\ref{sec:profileformat}, shown
	rather than described.

\end{description}
